\DocumentMetadata{
  pdfversion=2.0,
  pdfstandard=ua-2,
  lang=en-US,
  tagging=on,
  tagging-setup={math/setup=mathml-SE,tabsorder=structure}
}
\documentclass[11pt,letterpaper]{article}
\usepackage[margin=0.68in,includefoot]{geometry}
\usepackage{newcomputermodern}
\usepackage{amsmath,amscd,mathtools}
\usepackage{tikz,adjustbox}
\providecommand{\square}{\mdlgwhtsquare}
\usepackage[hidelinks]{hyperref}
\hypersetup{
  pdftitle={Algebra PhD Exam, May 1999},
  pdfauthor={University of Florida Department of Mathematics},
  pdfsubject={Graduate mathematics examination},
  pdfdisplaydoctitle=true,
  bookmarksopen=true
}
\setcounter{secnumdepth}{0}
\setlength{\parindent}{0pt}
\setlength{\parskip}{0.28em}
\setlength{\emergencystretch}{3em}
\setlength{\leftmargini}{2.15em}
\setlength{\leftmarginii}{2.65em}
\setlength{\leftmarginiii}{3em}
\pagestyle{empty}
\raggedbottom
\allowdisplaybreaks
\NewTaggingSocketPlug{block/list/label}{examproblem}{%
  \tagstructbegin{tag=Lbl}\tagstructend%
  \tagstructbegin{tag=\LBody}%
  \tagstructbegin{tag=\ProblemHeadingTag,title-o={\ProblemHeadingText}}%
  \tagmcbegin{tag=\ProblemHeadingTag,actualtext={\ProblemHeadingText}}%
  #2%
  \tagmcend%
  \tagstructend%
}
\newenvironment{examproblems}[2]{%
  \def\ProblemHeadingTag{#2}%
  \AssignTaggingSocketPlug{block/list/label}{examproblem}%
  \begin{enumerate}%
  \setcounter{enumi}{#1}%
  \renewcommand{\labelenumi}{\textbf{\theenumi.}}%
  \setlength{\topsep}{0.35em}%
  \setlength{\partopsep}{0pt}%
  \setlength{\itemsep}{0.48em}%
  \setlength{\parsep}{0.18em}%
}{\end{enumerate}}
\newenvironment{examsubparts}{%
  \AssignTaggingSocketPlug{block/list/label}{default}%
  \begin{enumerate}%
  \setlength{\topsep}{0.18em}%
  \setlength{\partopsep}{0pt}%
  \setlength{\itemsep}{0.16em}%
  \setlength{\parsep}{0.1em}%
}{\end{enumerate}}
\newenvironment{examinstructions}{%
  \par\begingroup\itshape
}{%
  \par\endgroup\vspace{0.18em}
}
\makeatletter
\renewcommand\subsection{\@startsection{subsection}{2}{0pt}%
  {-1.25ex plus -.25ex minus -.1ex}%
  {0.55ex plus .1ex}%
  {\normalfont\large\bfseries}}
\renewcommand{\@maketitle}{%
  \newpage\null\vskip -1em%
  {\footnotesize\bfseries
    \href{https://www.ufl.edu/}{University of Florida}, \href{https://math.ufl.edu/}{Department of Mathematics}\par}%
  \vskip -0.5em%
  \tagstructbegin{tag=H1,title={Algebra PhD Exam, May 1999}}%
  \tagpdfparaOff%
  \tagmcbegin{tag=H1}%
    {\LARGE\bfseries\href{https://gma.math.ufl.edu/exams/algebra-phd/}{Algebra PhD Exam}, May 1999\par}%
  \tagmcend%
  \tagpdfparaOn%
  \tagstructend%
  \vskip 0.25em%
  \tagmcbegin{artifact}%
    \hrule height 1pt%
  \tagmcend%
  \vskip 1em%
}
\makeatother
\title{Algebra PhD Exam, May 1999}
\author{}
\date{}
\begin{document}
\maketitle
\thispagestyle{empty}
\begin{examinstructions}
Please do 7 out of the 11 problems below.
\end{examinstructions}
\begin{examproblems}{0}{H2}
\def\ProblemHeadingText{Problem 1.}
\item
An algebraic extension field~\(F\) of~\(K\) is said to be normal over~\(K\) if every irreducible polynomial in \(K[x]\) that has a root in~\(F\) actually splits in \(F[x]\). Prove that an algebraic extension~\(F\) of~\(K\) is normal over~\(K\) if and only if for every irreducible \(f\in K[x]\), \(f\) factors in \(F[x]\) as a product of irreducible factors all of which have the same degree.
\def\ProblemHeadingText{Problem 2.}
\item
This problem has two parts.
\begin{examsubparts}
\item[\textnormal{(a).}]
Show that the additive group of rationals \(\mathbb Q\) is not free.
\item[\textnormal{(b).}]
Show that the group \(\mathbb Q^+\) of all positive rationals (under multiplication) is free abelian with basis \(\{p:p\text{ is prime in }\mathbb Z\}\).
\end{examsubparts}
\def\ProblemHeadingText{Problem 3.}
\item
Show that any simple group~\(G\) of order 60 is isomorphic to \(A_5\). (If you want, you may use the fact that for each \(n\geq 2\), \(A_n\) is the only subgroup of \(S_n\) of index 2.)
\def\ProblemHeadingText{Problem 4.}
\item
Let~\(R\) be an integral domain and for each maximal ideal~\(M\), consider the localization \(R_M\) as a subring of the quotient field of~\(R\). Show that \(\bigcap R_M=R\), where the intersection is taken over all maximal ideals~\(M\) of~\(R\).
\def\ProblemHeadingText{Problem 5.}
\item
This problem has two parts.
\begin{examsubparts}
\item[\textnormal{(a).}]
Define solvable group.
\item[\textnormal{(b).}]
Prove that any group of order 48 is solvable (without using Burnside’s \((p,q)\)-Theorem).
\end{examsubparts}
\def\ProblemHeadingText{Problem 6.}
\item
Let~\(R\) be a ring. Show that the following conditions on an~\(R\)-module~\(P\) are equivalent.
\begin{examsubparts}
\item[\textnormal{(i).}]
\(P\) is projective.
\item[\textnormal{(ii).}]
Every short exact sequence
\[
0\to A\xrightarrow{f}B\xrightarrow{g}P\to 0
\]
 is split exact.
\item[\textnormal{(iii).}]
There is a free module~\(F\) and an~\(R\)-module~\(K\) such that \(F\cong K\oplus P\).
\end{examsubparts}
\def\ProblemHeadingText{Problem 7.}
\item
Let~\(C\) and~\(D\) be categories, and let~\(S\) and~\(T\) be covariant functors from~\(C\) to~\(D\).
\begin{examsubparts}
\item[\textnormal{(a).}]
Define a natural isomorphism \(\alpha:S\to T\).
\item[\textnormal{(b).}]
If~\(B\) is a unitary left module over a ring~\(R\) with identity, show that there is a “natural” isomorphism of modules
\[
\alpha_B:R\otimes_R B\cong B.
\]
\end{examsubparts}
\def\ProblemHeadingText{Problem 8.}
\item
Let~\(R\) be a commutative ring with identity.
\begin{examsubparts}
\item[\textnormal{(a).}]
Let \(\mathcal A\) be an ideal in~\(R\), and assume that~\(M\) is a finitely generated~\(R\)-module such that \(\mathcal A.M=M\). Show that there is some \(a\in\mathcal A\) satisfying \((1+a)M=0\).
\item[\textnormal{(b).}]
Let~\(M\) be a finitely-generated~\(R\)-module. Show that \(J(R).M=M\) implies \(M=\{0\}\). (Here, \(J(R)\) denotes the Jacobson radical of~\(R\).)
\end{examsubparts}
\def\ProblemHeadingText{Problem 9.}
\item
This problem has two parts.
\begin{examsubparts}
\item[\textnormal{(a).}]
Show that if~\(R\) is a unique factorization domain, then~\(R\) is integrally closed.
\item[\textnormal{(b).}]
Find the integral closure of \(\mathbb C[x^5,x^7]\).
\end{examsubparts}
\def\ProblemHeadingText{Problem 10.}
\item
This problem has two parts.
\begin{examsubparts}
\item[\textnormal{(a).}]
Show that \(\mathbb C[x,y]/(xy-1)\) (quotient of polynomial ring in 2 variables) is not isomorphic to \(\mathbb C[t]\) (polynomial ring in 1 variable).
\item[\textnormal{(b).}]
Show that the set
\[
\{(m,n):m,n\in\mathbb Z\}
\]
 is not an algebraic variety in \(\mathbb C^2\).
\end{examsubparts}
\def\ProblemHeadingText{Problem 11.}
\item
This problem has three parts.
\begin{examsubparts}
\item[\textnormal{(a).}]
Find the Galois group of \(K:\mathbb Q\), where~\(K\) is a splitting field over \(\mathbb Q\) for \(t^4+t^2-6\).
\item[\textnormal{(b).}]
Write down the subgroups of the Galois group.
\item[\textnormal{(c).}]
Write down the corresponding fixed fields.
\end{examsubparts}
\end{examproblems}
\end{document}
